\documentclass{article}
\usepackage{apstats-poster}
% Formula IDs: phat-mean, phat-sd, xbar-mean, xbar-sd, ten-pct-condition
\colorlet{PosterAccent}{UnitBlue}

\begin{document}
\begin{posterpage}
\PosterHeader[70pt]{UnitBlue}{2}{PROBABILITY, RANDOM VARIABLES \& SAMPLING DISTRIBUTIONS}
  {P08}{SAMPLING DISTRIBUTIONS \& THE CLT}
  {new 2.12, 3.1--3.2, 4.1 \quad\textbullet\quad old 5.1--5.7}

\PosterZone{4.72in}{15.90in}{ZONE A · THE FORMULAS + CONDITIONS}{%
  \MiniHeading{1 · RANDOM: START WITH A RANDOM SAMPLE}\par\vspace{.18in}
  \begin{minipage}[t]{9.80in}
    \TextCard{\centering\fontsize{33pt}{40pt}\selectfont
      \textbf{SAMPLE PROPORTION \(\hat p\)}\par
      Population proportion \(p\)}
    \vspace{.16in}
    \FormulaCardSized{57pt}{\mu_{\hat p}=p}
    \vspace{.16in}
    \FormulaCardSized{53pt}{\sigma_{\hat p}=\sqrt{\frac{p(1-p)}n}}
  \end{minipage}\hfill
  \begin{minipage}[t]{9.80in}
    \TextCard{\centering\fontsize{33pt}{40pt}\selectfont
      \textbf{SAMPLE MEAN \(\bar x\)}\par
      Population mean \(\mu\), SD \(\sigma\)}
    \vspace{.16in}
    \FormulaCardSized{57pt}{\mu_{\bar x}=\mu}
    \vspace{.16in}
    \FormulaCardSized{53pt}{\sigma_{\bar x}=\frac{\sigma}{\sqrt n}}
  \end{minipage}
  \vspace{.24in}
  \TextCard{\centering\fontsize{31pt}{39pt}\selectfont
    \textbf{2 · 10\%:} \(\boldsymbol{n<0.10\cdot N}\) when sampling without replacement.\par
    Draws are close enough to independent to use these SD formulas.}
  \vspace{.24in}
  \begin{minipage}[t]{9.80in}
    \MiniHeading{3 · LARGE COUNTS}\par\vspace{.13in}
    {\fontsize{31pt}{40pt}\selectfont
      The distribution of \(\hat p\) is approximately normal when\par\vspace{.10in}
      \(\boldsymbol{np\ge10\quad\text{and}\quad n(1-p)\ge10}\).}
  \end{minipage}\hfill
  \begin{minipage}[t]{9.80in}
    \RaggedRight
    \MiniHeading{3 · NORMAL / CENTRAL LIMIT THEOREM}\par\vspace{.13in}
    {\fontsize{29pt}{37pt}\selectfont
      Normal population \(\Rightarrow\) normal \(\bar x\), for any \(n\).\par
      Otherwise, CLT gives an \textbf{approximate} normal shape for large \(n\).
      The \(n\ge30\) guideline is approximate; strong skew can need a larger sample.}
  \end{minipage}
}

\PosterZone{16.05in}{20.70in}{ZONE B · PARAMETERS VS STATISTICS}{%
  \begin{minipage}[t]{8.0in}
    \vspace{0pt}
    {\fontsize{31pt}{39pt}\selectfont
      \renewcommand{\arraystretch}{1.24}
      \begin{tabular}{lcc}
        & \textbf{Parameter} & \textbf{Statistic}\\\midrule
        Mean & \(\mu\) & \(\bar x\)\\
        SD & \(\sigma\) & \(s\)\\
        Proportion & \(p\) & \(\hat p\)\\
        Slope & \(\beta\) & \(b\)
      \end{tabular}}
  \end{minipage}\hfill
  \begin{minipage}[t]{11.65in}
    \vspace{0pt}\RaggedRight
    {\fontsize{28pt}{36pt}\selectfont
      \textbf{Parameters are Greek or unhatted. Statistics wear a hat or a bar}
      (or use the sample symbols \(s\) and \(b\)).\par\vspace{.16in}
      \textbf{Unbiased:} the sampling distribution is centered at the parameter.\par\vspace{.16in}
      Variability shrinks with \(\sqrt n\): \textbf{quadrupling \(n\) halves the SD}.}
  \end{minipage}
}

\PosterZone{20.85in}{23.85in}{ZONE C · SAY IT IN WORDS}{%
  \SentenceFrame{“The sampling distribution of \PosterBlank{2.8in} is the
    distribution of that statistic from all possible samples of size
    \(n=\) \PosterBlank{1.0in} from \PosterBlank{3.1in}.”}
  \vspace{.12in}
  {\fontsize{26pt}{33pt}\selectfont
    One sample gives one statistic; repeated samples give a distribution.}
}

\PosterZone{24.00in}{27.45in}{ZONE D · WATCH OUT}{%
  \begin{minipage}[t]{9.65in}
    \RaggedRight
    \MiniHeading{KEEP \(n\) AND \(N\) SEPARATE}\par\vspace{.12in}
    {\fontsize{30pt}{38pt}\selectfont
      \(n\) is sample size; \(N\) is population size.\par
      Sample size controls these SDs.\par
      Use \(N\) for the 10\% check.}
  \end{minipage}\hfill
  \begin{minipage}[t]{9.95in}
    \RaggedRight
    \MiniHeading{SD VS ESTIMATED SD}\par\vspace{.12in}
    {\fontsize{29pt}{37pt}\selectfont
      \(\sigma_{\bar x}=\sigma/\sqrt n\) needs population SD \(\sigma\).
      With sample SD \(s\), \(s/\sqrt n\) is an \textbf{SE} estimating
      the SD of \(\bar x\). See P13.}
  \end{minipage}
}
\WorksheetTag{u5\_l1--7}
\end{posterpage}
\end{document}
